\documentclass[12pt,a4paper]{article}
\usepackage[utf8]{inputenc}
\usepackage[T1]{fontenc}
\usepackage{lmodern}
\usepackage{amsmath}
\usepackage[spanish]{babel}
\usepackage{microtype}
\usepackage{hyperref}
\usepackage{bookmark}
\usepackage{graphicx}
\usepackage{geometry}
\geometry{a4paper, margin=1in}
\emergencystretch=2.5em
\renewcommand{\_}{\textunderscore\allowbreak}

\title{Instructivo de Uso: Credit Manager}
\author{Documentación Oficial}
\date{\today}

\begin{document}

\maketitle
\newpage

\tableofcontents
\newpage

\section{Introducción}
Bienvenido a \textbf{Credit Manager}, la plataforma integral diseñada para simplificar y agilizar la administración diaria de préstamos, cobranzas, cheques, inversores y finanzas en tu empresa. 

Este instructivo fue elaborado pensando en las personas que utilizan el sistema en su labor cotidiana (operadores comerciales, analistas de crédito, cajeros, responsables de cobranzas y personal de tesorería). Aquí encontrarás explicaciones claras, sin tecnicismos innecesarios, y guías prácticas paso a paso para que puedas operar cada módulo con total confianza y eficacia.

\newpage

\section{¿Qué Podés Hacer con Credit Manager?}
El sistema reúne en un solo lugar todas las herramientas necesarias para la operatoria financiera de tu compañía:

\begin{itemize}
    \item \textbf{Simular y Otorgar Créditos al Instante:} Calculá cronogramas de cuotas en el momento frente al cliente bajo el Sistema Francés, discriminando con exactitud el capital, los intereses y el IVA.
    \item \textbf{Ficha Completa de Clientes:} Consultá en una única pantalla los datos personales, historial de pagos, situación crediticia en la Central de Deudores del BCRA y control de listas de seguridad (RePET).
    \item \textbf{Cobranzas Ágiles (Uno a Uno o por Planilla):} Registrá cobros individuales en caja o importá directamente el archivo de Excel que envía el banco o el ente pagador para imputar cientos de cuotas de una sola vez.
    \item \textbf{Facturación Electrónica Automática:} Emití comprobantes con autorización electrónica (CAE) conectándote de forma directa con ARCA (ex AFIP) con un solo clic.
    \item \textbf{Control Total de Cheques Físicos y Electrónicos (E-Cheqs):} Llevá la custodia de valores, calculá descuentos financieros en tiempo real y conciliá depósitos o débitos con tus cuentas bancarias.
    \item \textbf{Gestión de Inversores y Fondeo:} Estructurá series de pagarés o préstamos privados, liquidá intereses mensuales y gestioná renovaciones o rescates de capital con total transparencia.
    \item \textbf{Tesorería, Bancos y Posiciones Fiscales:} Importá resúmenes bancarios para conciliar movimientos, controlá facturas de proveedores y liquidá tus posiciones mensuales de IVA e Ingresos Brutos (IIBB).
    \item \textbf{Papelería Legal Automática:} Descargá pagarés, contratos de cesión y solicitudes completadas automáticamente con los datos del cliente en formatos Word y PDF listos para firmar.
\end{itemize}

\newpage

\section{¿Cómo Funciona el Sistema?}
Para que no tengas que preocuparte por cuestiones técnicas, \textbf{Credit Manager} está organizado en tres partes sencillas:
\begin{enumerate}
    \item \textbf{La Pantalla Web (lo que ves y usás):} Es el panel visual interactivo al que ingresás desde tu navegador web habitual (Google Chrome, Microsoft Edge, etc.). Es rápido, intuitivo y podés usarlo desde cualquier computadora autorizada.
    \item \textbf{El Motor de Procesamiento (el cerebro del sistema):} Trabaja en segundo plano ejecutando de forma inmediata todos los cálculos matemáticos: liquidación de cuotas, devengamiento de intereses, validaciones fiscales y cruces de seguridad.
    \item \textbf{La Base de Datos Segura (el archivo digital):} Guarda toda la información de clientes, préstamos, cobranzas y comprobantes con resguardo permanente, impidiendo que los datos se borren por error o se dupliquen.
\end{enumerate}

\newpage

\section{Puesta en Marcha y Configuración Inicial}
\subsection{Datos de la Empresa}
Antes de comenzar a operar, el administrador del sistema configura los datos básicos de tu organización en el archivo de entorno (\texttt{.env}) o desde el panel de configuración:
\begin{itemize}
    \item \textbf{Razón Social:} El nombre legal de tu compañía, que aparecerá en los recibos, contratos y reportes.
    \item \textbf{CUIT:} La Clave Única de Identificación Tributaria para la facturación electrónica ante ARCA/AFIP.
    \item \textbf{Día de Corte:} El día habitual del mes fijado para el cálculo de vencimientos y liquidaciones de cartera (por ejemplo, el día 28 de cada mes).
\end{itemize}

\subsection{Inicio de la Aplicación}
Para poner en marcha el sistema en el servidor de la empresa, el equipo técnico o administrador puede iniciarlo mediante contenedores (Docker) o mediante el servidor local. Una vez iniciado, todos los usuarios pueden ingresar abriendo el navegador web en la dirección de red local asignada (por ejemplo, \texttt{http://localhost:5173} o la dirección interna de la empresa).
\newpage

\section{Uso y Comportamiento del Sistema}

\subsection{Dashboard (Panel Principal)}
El \textbf{Dashboard} es la primera pantalla que ves al ingresar a Credit Manager. Funciona como el tablero de control de tu empresa: te muestra en tiempo real, mediante tarjetas y gráficos claros, cuánta plata hay prestada en la calle, cuánto se cobró en el mes, qué porcentaje de la cartera está al día y cuánta deuda se encuentra en mora.

\textbf{¿Qué información encontrás en este panel?}
\begin{itemize}
    \item \textbf{Tarjetas de Resumen (KPIs):} Te muestran de un vistazo el Capital Total Activo, el total de dinero recaudado en el mes corriente, el Saldo en Mora y la cantidad de créditos vigentes.
    \item \textbf{Evolución de Cobranzas vs. Vencimientos:} Un gráfico que compara las cuotas que tenían que vencer este mes contra lo que efectivamente pagaron los clientes.
    \item \textbf{Distribución de Mora:} Un gráfico circular que te muestra visualmente qué porcentaje de la cartera está al día y cuánto se reparte en atrasos leves (menos de 30 días), medios o graves.
    \item \textbf{Proyección de Fondos:} Una estimación de cuánta plata va a ingresar a la empresa en los próximos meses si los clientes pagan según su cronograma habitual.
\end{itemize}

\textbf{Paso a paso: Cómo filtrar y consultar la información del Dashboard}
\begin{enumerate}
    \item En la barra de filtros superior, elegí la \textbf{Fecha de Corte} (por defecto toma el día de hoy, pero podés elegir el último día del mes pasado si querés analizar el cierre anterior).
    \item Si tu empresa administra carteras de distintos socios o inversores, usá los desplegables de \textbf{Dueños} u \textbf{Originadores} para ver los números de un socio en particular (si los dejás en blanco, el sistema te muestra el total consolidado de la empresa).
    \item Hacé clic en el botón azul \textbf{"Aplicar Filtros"}. Todos los números, barras y gráficos de la pantalla se actualizarán en un instante.
    \item Si necesitás presentar estos números en una reunión o enviárselos a gerencia, hacé clic en \textbf{"Exportar a PDF"} para descargar un informe prolijo y listo para imprimir.
\end{enumerate}

\subsection{Cartera de Créditos}

\subsubsection{Simular Cuotas}
Esta herramienta te permite armar presupuestos y calcular planes de pago en el momento, ya sea con el cliente en el mostrador o mientras hablás por teléfono con un asesor comercial. Utiliza el tradicional \textbf{Sistema Francés} (cuotas fijas y consecutivas).

\textbf{Paso a paso: Cómo simular un préstamo para un cliente}
\begin{enumerate}
    \item Ingresá a \textbf{Cartera de Créditos} y hacé clic en \textbf{Simular Cuotas}.
    \item En el campo \textbf{Capital a financiar}, escribí el monto de dinero que el cliente solicita (por ejemplo, \$150.000).
    \item En \textbf{Plazo (Cuotas)}, indicá en cuántos meses lo va a devolver (por ejemplo, 12, 18 o 24 cuotas).
    \item En \textbf{TNA c/IVA (\%)}, ingresá la Tasa Nominal Anual que acordaste para ese tipo de crédito.
    \item Indicá la \textbf{Fecha de Emisión} (cuándo se entregaría la plata) y el \textbf{Día de Vencimiento} en el que el cliente deberá pagar cada mes (por ejemplo, los días 10).
    \item Si la empresa le otorga un tiempo de gracia antes de empezar a pagar, indicalo en el campo \textbf{Meses de Gracia} (si empieza a pagar al mes siguiente, dejalo en 0).
    \item Hacé clic en el botón \textbf{"Simular Cuotas"}.
    \item Inmediatamente aparecerá en pantalla el cuadro completo de cuotas, mostrándote para cada mes el valor final de la cuota y su desglose exacto: cuánto va a devolución de capital, cuánto a intereses y cuánto a IVA.
    \item Si el cliente quiere llevarse la propuesta impresa o querés enviársela por WhatsApp, hacé clic en \textbf{"Exportar a Excel"}.
\end{enumerate}

\subsubsection{Clientes}
En este módulo administrás el padrón completo de personas y empresas a las que tu organización les otorga financiamiento.

\textbf{Ficha Integral del Cliente (Dashboard de Cliente):}
Es como el legajo digital del cliente. Al escribir su DNI, CUIL o apellido en el buscador, el sistema abre su ficha completa donde podés ver:
\begin{itemize}
    \item Sus datos personales, teléfonos y domicilio verificado.
    \item Su calificación crediticia oficial en la Central de Deudores del Banco Central (BCRA).
    \item Si registra alertas de prevención de lavado o terrorismo en las listas oficiales (RePET).
    \item Todos los préstamos que tuvo y tiene con tu empresa, con el detalle de cuotas pagadas y atrasos.
    \item Su cuenta corriente histórica para saber si tiene saldo a favor o deudas pendientes.
\end{itemize}

\textbf{Paso a paso: Cómo dar de alta a un nuevo cliente}
\begin{enumerate}
    \item En el menú lateral, andá a \textbf{Cartera de Créditos} y seleccioná \textbf{Clientes}.
    \item Hacé clic en el botón superior \textbf{"+ Nuevo Cliente"}.
    \item Completá los datos identificatorios:
          \begin{itemize}
              \item \textbf{CUIL/CUIT:} 11 dígitos numéricos sin guiones ni espacios.
              \item \textbf{DNI, Apellido y Nombre completo.}
              \item \textbf{Sexo y Fecha de Nacimiento.}
          \end{itemize}
    \item Cargá sus vías de contacto: Teléfono celular principal, un teléfono alternativo y su correo electrónico para el envío de avisos de vencimiento.
    \item Completá su \textbf{Domicilio}: Calle, Número, Piso/Depto, Código Postal, Localidad y Provincia.
    \item Ingresá sus \textbf{Datos Laborales y Bancarios}: Empleador donde trabaja, sueldo neto mensual y su CBU o número de cuenta bancaria donde suele cobrar.
    \item Hacé clic en \textbf{"Guardar Cliente"}. El sistema validará que el CUIL no esté repetido y revisará automáticamente las listas de seguridad para garantizar que la persona esté habilitada para operar.
\end{enumerate}

\textbf{Acciones útiles en la lista de clientes:}
Al lado de cada cliente tenés botones directos para: ver su legajo de papeles digitalizados (\textit{Ver Detalles}), consultar su saldo actual (\textit{Ver Cuenta Corriente}) o saltar directamente a ver sus créditos activos (\textit{Créditos}).

\subsubsection{Créditos}
Este es el corazón operativo del sistema: desde acá otorgás nuevos préstamos, supervisás la cartera activa y controlás qué créditos están al día y cuáles tienen atrasos.

\textbf{Estados de un crédito (¿qué significan los colores?):}
\begin{itemize}
    \item \textbf{Al Día (Verde):} El cliente pagó todas las cuotas vencidas o aún no llegó el primer vencimiento.
    \item \textbf{Mora Leve (Amarillo):} Registra cuotas impagas con atrasos de entre 1 y 30 días.
    \item \textbf{Mora Media (Naranja):} Registra atrasos de entre 31 y 90 días.
    \item \textbf{Mora Grave (Rojo):} Atraso superior a 90 días que requiere gestión intensiva de cobranzas.
    \item \textbf{Cancelado (Gris / Azul):} El préstamo fue saldado por completo y no debe nada.
\end{itemize}

\textbf{Paso a paso: Cómo otorgar un crédito nuevo (Asistente en 5 pasos)}
Cuando un cliente solicita un préstamo, hacé clic en el botón verde \textbf{"+ Nuevo Crédito"} para abrir el asistente guiado:
\begin{enumerate}
    \item \textbf{Paso 1 - Buscar al Cliente:} Escribí el DNI o CUIL del solicitante. Si ya es cliente de la empresa, el sistema completará sus datos automáticamente. Si es una persona nueva, podés cargarla en el momento sin salir del asistente.
    \item \textbf{Paso 2 - Verificar Requisitos y Documentación:} Revisá que los datos laborales estén vigentes y confirmá que la persona cumpla con las políticas crediticias de la empresa (edad mínima y máxima, antigüedad laboral requerida, etc.).
    \item \textbf{Paso 3 - Evaluación de Riesgo y Capacidad de Pago:} El sistema consulta en línea la situación crediticia en el BCRA y analiza el sueldo del cliente para calcular la \textbf{Cuota Afectable Máxima} (asegurando que el valor de la cuota no supere el porcentaje permitido de sus ingresos).
    \item \textbf{Paso 4 - Condiciones Financieras:} Ingresá el monto de dinero neto que se le va a entregar en mano al cliente y elegí el plan de tasas y comisiones correspondiente al socio originador. El sistema calculará automáticamente los gastos administrativos, sellados provinciales y el capital bruto del préstamo.
    \item \textbf{Paso 5 - Cronograma, Transferencias y Confirmación:} Vas a ver en pantalla la tabla definitiva de cuotas calculadas bajo el Sistema Francés. Indicá el CBU o cuenta bancaria a la que se le transferirá el dinero al cliente y, si corresponde, las comisiones a los intermediarios. Podés previsualizar la documentación en PDF y, al hacer clic en \textbf{"Confirmar y Otorgar"}, el crédito queda formalmente asentado y vigente en el sistema.
\end{enumerate}

\textbf{Herramientas para créditos existentes en la grilla:}
En cada fila de crédito contás con accesos directos muy útiles:
\begin{itemize}
    \item \textbf{Ver Cuenta Corriente:} Te abre una ventana donde ves cuota por cuota qué día venció, cuánto pagó, cuánto debe y qué recibos tiene aplicados.
    \item \textbf{Ver Transferencias:} Te muestra el detalle del desembolso bancario: fecha de transferencia, monto exacto, CBU de destino y el comprobante bancario adjunto.
    \item \textbf{Ver Legajo Digital:} Te permite consultar y descargar en pantalla las fotos o PDFs del DNI del cliente, sus recibos de sueldo y los contratos firmados.
\end{itemize}

\subsubsection{Cobranzas}
En este módulo se asientan todos los pagos que ingresan a la empresa. 

\textbf{¿Cómo imputa el dinero el sistema? (Cobranza en cascada)}
Para evitarle problemas y sobrecargos al cliente, el sistema aplica siempre una regla de cobro justa y transparente:
\begin{itemize}
    \item \textbf{Primero lo más viejo:} El pago cancela siempre la cuota más antigua que esté impaga, evitando que sigan corriendo intereses punitorios.
    \item \textbf{Orden financiero dentro de la cuota:} Cada peso que entra cancela primero los impuestos (IVA), luego los intereses devengados y finalmente el capital del préstamo.
\end{itemize}

\textbf{Paso a paso: Cómo registrar un cobro individual (en caja o transferencia)}
\begin{enumerate}
    \item Andá a \textbf{Cartera de Créditos} > \textbf{Cobranzas} y hacé clic en \textbf{"Nueva Cobranza"}.
    \item Buscá la deuda escribiendo el DNI, CUIL o número de crédito del cliente.
    \item Ingresá el \textbf{Monto Cobrado} (el dinero exacto que pagó) y la \textbf{Fecha de Pago}.
    \item Elegí el \textbf{Tipo de Cobranza}:
          \begin{itemize}
              \item \textit{Común / Cuota:} Si el cliente está pagando su cuota mensual habitual o entregando dinero a cuenta.
              \item \textit{Cancelación Total Anticipada:} Si el cliente viene a cancelar toda su deuda pendiente de golpe. En este caso, el sistema descuenta automáticamente los intereses futuros no devengados para cobrarle únicamente el capital real que debe.
          \end{itemize}
    \item Hacé clic en \textbf{"Procesar Cobranza"}. El sistema saldará la cuota, actualizará el estado del crédito y te mostrará el recibo de pago listo para entregarle al cliente.
\end{enumerate}

\textbf{Paso a paso: Cómo procesar una cobranza masiva con Excel}
Cuando cobrás a través de débitos bancarios, descuentos por recibo de sueldo o redes de cobranza, te envían un archivo de Excel con cientos de pagos juntos. Para no cargarlos uno por uno a mano:
\begin{enumerate}
    \item Ingresá a \textbf{Cobranzas} y seleccioná la solapa o botón \textbf{"Cobranza Masiva"}.
    \item Hacé clic en \textbf{"Seleccionar Archivo"} y elegí la planilla de Excel (\texttt{.xlsx}) en tu computadora.
    \item En la pantalla vas a ver un selector muy simple: indicá en qué columna de tu Excel está el identificador del cliente (por ejemplo, Columna A: CUIL o DNI) y en qué columna está el dinero pagado (por ejemplo, Columna B: Importe).
    \item Hacé clic en \textbf{"Validar Planilla"}. El sistema revisará el archivo sin guardar nada todavía y te mostrará un resumen claro: cuántos registros coinciden perfectamente con deudas activas y si hay algún DNI que no se encontró en la base.
    \item Si todo coincide, hacé clic en \textbf{"Confirmar e Imputar Cobranzas"}. En unos pocos segundos, el sistema cancelará todas las cuotas del archivo de forma automática.
\end{enumerate}

\textbf{¿Qué pasa si me equivoqué de archivo o cargué un cobro mal? (Reversión segura)}
No te preocupes: el sistema está protegido contra errores humanos.
\begin{itemize}
    \item Si te equivocaste en un cobro individual, podés buscarlo en la grilla y hacer clic en el botón de eliminar cobro para anularlo.
    \item Si subiste una planilla masiva equivocada, andá a la solapa \textbf{Procesos}, buscá el lote que acabás de procesar (está identificado con fecha y hora) y hacé clic en \textbf{"Revertir Proceso"}. El sistema anulará el lote entero de un solo golpe y devolverá todas las cuotas a su estado anterior, como si nunca hubieras subido la planilla.
\end{itemize}

\subsubsection{Facturación (Conexión con ARCA / AFIP)}
Este módulo se encarga de emitir las facturas electrónicas oficiales ante la agencia tributaria (ARCA, ex AFIP) por las cobranzas que realizaste en el mes.

\textbf{¿Cómo funciona impositivamente?}
Cuando un cliente paga su cuota, la porción correspondiente a la devolución del capital no paga impuestos (concepto No Gravado). En cambio, los intereses cobrados sí llevan IVA y deben facturarse formalmente. El sistema hace esta separación de forma automática y matemática para que vos no tengas que calcular nada a mano.

\textbf{Paso a paso: Cómo facturar las cobranzas pendientes del mes}
\begin{enumerate}
    \item Ingresá a \textbf{Cartera de Créditos} > \textbf{Facturación}.
    \item En la solapa \textbf{Cobranzas Pendientes de Facturar}, vas a ver el listado de todos los cobros que aún no tienen comprobante fiscal. Podés usar los filtros si querés facturar solo los cobros de un socio originador o un rango de fechas determinado.
    \item Seleccioná las cobranzas que querés facturar (podés tildar la casilla superior para marcar todas juntas).
    \item En el campo \textbf{Fecha de Emisión}, elegí la fecha de tus facturas (el sistema te avisa si la fecha respeta los plazos legales de ARCA para que nunca tengas un rechazo).
    \item Hacé clic en el botón verde \textbf{"Generar Facturas Electrónicas"}.
    \item El sistema se conecta en el acto con los servidores de ARCA/AFIP, solicita el código de autorización (CAE) de cada factura, le asigna su punto de venta y número oficial correlativo, y deja los cobros registrados como facturados.
\end{enumerate}

\textbf{Descargas para el contador y los clientes:}
\begin{itemize}
    \item \textbf{Descargar Libro IVA Ventas (Excel):} Te descarga la planilla oficial con fecha, número de factura, CUIT del receptor, importes gravados, débito de IVA y número de CAE, lista para enviarle a tu contador a fin de mes.
    \item \textbf{Descargar Facturas (ZIP):} Te genera un archivo comprimido con todos los comprobantes individuales en formato PDF (con el logo de tu empresa, el detalle legal y el código de barras de ARCA) listos para imprimir o mandar por mail a los clientes.
\end{itemize}

\subsubsection{Carteras (Compra y Venta de Carteras)}
Este módulo te permite comprar o vender paquetes de préstamos en bloque a otras instituciones financieras, bancos o inversores.

\textbf{¿Qué significa vender una cartera (Cesión de Créditos)?}
Significa que tu empresa transfiere los derechos de cobro de un grupo de préstamos a un tercero para hacerse de dinero en efectivo hoy en lugar de esperar a cobrar cuota por cuota mes a mes.
\begin{itemize}
    \item \textbf{Con Recurso:} Si algún cliente de la cartera vendida deja de pagar, tu empresa sigue siendo responsable frente al comprador y debe compensar esa cuota impaga.
    \item \textbf{Sin Recurso:} El comprador asume todo el riesgo crediticio; si el cliente no paga, la pérdida es del comprador y tu empresa no debe nada.
\end{itemize}

\textbf{Paso a paso: Cómo armar y cotizar una Venta de Cartera}
\begin{enumerate}
    \item Andá a \textbf{Carteras} y hacé clic en \textbf{"Nueva Operación"} seleccionando \textbf{Modalidad Venta}.
    \item Elegí el \textbf{Socio Comercial} que te va a comprar los créditos.
    \item Indicá la \textbf{Fecha de Corte} de la operación y la \textbf{Tasa Nominal Anual (TNA) de Descuento} que acordaron con la contraparte.
    \item Indicá si la venta se realiza \textbf{Con Recurso} o \textbf{Sin Recurso}.
    \item Usá los filtros para seleccionar qué créditos entran en el paquete: podés filtrar por fecha de otorgamiento, por fecha de vencimiento de las cuotas o excluir a clientes que tengan mora.
    \item En la pantalla de previsualización vas a ver el cálculo del \textbf{Valor Actual (VA)}: el sistema descuenta matemáticamente cada cuota futura y te dice con precisión cuánta plata en efectivo te van a pagar hoy por todo el paquete. Podés desmarcar créditos puntuales si preferís dejarlos afuera.
    \item Hacé clic en \textbf{"Confirmar Venta"}.
    \item Una vez confirmada, desde la grilla de operaciones podés descargar con un solo clic:
          \begin{itemize}
              \item El \textbf{Contrato Formal de Cesión} en formato Word o PDF listo para imprimir y firmar ante escribano.
              \item El \textbf{Excel de la Operación} con las solapas de Créditos, Cuotas y Flujo mensual de cobros.
              \item El \textbf{Paquete de Legajos Digitales (ZIP)} con las fotos de los DNI y recibos de sueldo de todos los clientes de esa cartera.
          \end{itemize}
\end{enumerate}

\textbf{Paso a paso: Cómo hacer la liquidación de cobranzas de carteras cedidas}
Si vendiste una cartera pero tu empresa sigue encargándose de cobrar las cuotas, tenés que rendirle cuentas periódicamente al comprador:
\begin{enumerate}
    \item Entrá a la solapa \textbf{Liquidaciones de Cartera}.
    \item Hacé clic en \textbf{"Nueva Liquidación"}.
    \item Elegí el socio comercial comprador y el período de cobro que querés liquidar.
    \item El sistema agrupa automáticamente todas las cuotas que tus clientes pagaron durante ese período y calcula el importe exacto a transferirle al socio.
    \item Hacé clic en \textbf{"Confirmar Liquidación"} y registrá la fecha y comprobante de la transferencia bancaria realizada.
\end{enumerate}

\subsubsection{Reportes de Créditos}
Este submódulo reúne todos los informes analíticos y las declaraciones obligatorias que te solicitan los organismos oficiales.

\textbf{1. Reporte de Saldos de Cartera (¿Cuánta plata tenemos en la calle?)}
Te permite responder en segundos cualquier consulta sobre el estado patrimonial de tus créditos a cualquier fecha:
\begin{itemize}
    \item Podés elegir si querés ver solo los préstamos propios, solo los vendidos a terceros o ambos consolidados.
    \item Podés agrupar los datos como te sea más cómodo: por cliente, por socio originador, por mes de vencimiento o por número de crédito.
    \item En pantalla y en la descarga de Excel vas a ver el desglose exacto de: Capital vivo adeudado, Intereses devengados, IVA y el Saldo Total exigible.
\end{itemize}

\textbf{2. Reporte Oficial para el BCRA (Proveedores No Financieros de Crédito - PNFC)}
Si tu empresa otorga préstamos, tiene la obligación legal de informar mensualmente el estado de sus deudores al Banco Central de la República Argentina:
\begin{enumerate}
    \item En la solapa de \textbf{Reportes Regulados}, elegí la opción \textbf{Régimen Informativo BCRA}.
    \item Seleccioná la fecha de corte del mes que querés presentar.
    \item El sistema clasifica automáticamente a cada cliente según los días de mora que tenga:
          \begin{itemize}
              \item \textit{Situación 1 (Normal):} Al día o atraso menor a 31 días.
              \item \textit{Situación 2 (Riesgo Bajo):} Atraso de 32 a 90 días.
              \item \textit{Situación 3 (Riesgo Medio):} Atraso de 91 a 180 días.
              \item \textit{Situación 4 (Riesgo Alto):} Atraso de 181 a 365 días.
              \item \textit{Situación 5 (Irrecuperable):} Atraso mayor a 365 días.
          \end{itemize}
    \item Hacé clic en el botón \textbf{"Descargar Archivos Oficiales (.ZIP)"}.
    \item El sistema te descarga un archivo comprimido con los tres archivos de texto obligatorios (\texttt{PROVEEDORES.TXT}, \texttt{IMPORTES.TXT} y \texttt{TASA.TXT}), formateados con los anchos de columna y caracteres exactos que exige el validador oficial del BCRA. ¡Listos para subir sin rechazos!
    \item También podés descargar la planilla de control en Excel para auditar los números antes de la presentación.
\end{enumerate}

\textbf{3. Reporte para la UIF (Prevención de Lavado de Activos y PLA/FT)}
Te permite generar en segundos la auditoría requerida por la Unidad de Información Financiera:
\begin{itemize}
    \item \textbf{Padrón de Clientes:} Nómina completa con DNI, CUIL, capital total tomado y saldo de deuda actual.
    \item \textbf{Matriz de Riesgo:} Detalle del puntaje de riesgo asignado a cada tomador según sus antecedentes.
    \item \textbf{Registro de Cruce RePET:} Constancia fechada de que cada cliente fue cotejado contra la base pública de personas vinculadas al terrorismo.
\end{itemize}

\subsection{Cartera de Cheques}

El módulo de \textbf{Cartera de Cheques} funciona como la caja fuerte digital de tu empresa: te permite llevar el control exacto de todos los cheques que tenés en mano, tanto los tradicionales de papel como los cheques electrónicos (\textbf{E-Cheqs}). 

El sistema distingue con total claridad entre:
\begin{itemize}
    \item \textbf{Cheques Propios:} Son los cheques emitidos por tu propia empresa para pagar a proveedores o socios.
    \item \textbf{Cheques de Terceros:} Son los valores que recibís de clientes como pago de sus cuotas o que comprás en el mercado financiero mediante descuento de cheques.
\end{itemize}

\subsubsection{Operar Cheques}
Desde esta pantalla administrás el ingreso de cheques, calculás el dinero en mano al descontarlos, controlás quién tiene la tenencia de cada valor y conciliás los movimientos con tus cuentas bancarias.

\textbf{Padrón de Emisores y Operadores:}
Cada persona o empresa que emite o compra cheques se registra con su CUIT (11 dígitos sin guiones), Razón Social, Teléfono y Calificación Crediticia (\textit{Excelente}, \textit{Bueno}, \textit{Regular}, \textit{Malo} o \textit{Rechazado}). Al iniciar el sistema, tu propia empresa ya viene cargada automáticamente con calificación Excelente.

\textbf{Paso a paso: Cómo dar de alta un cheque en cartera}
\begin{enumerate}
    \item En el menú lateral, ingresá a \textbf{Cartera de Cheques} y hacé clic en el botón verde \textbf{"+ Nuevo Cheque"}.
    \item Elegí el \textbf{Emisor (Librador)}:
          \begin{itemize}
              \item Si el cheque lo emitió tu propia empresa, elegí tu razón social y el sistema lo marcará automáticamente con la etiqueta \texttt{PROPIO}.
              \item Si el cheque te lo entregó un cliente o tercero, elegí su nombre de la lista (o hacé clic en \texttt{+ Nuevo} si es la primera vez que opera).
          \end{itemize}
    \item Seleccioná el \textbf{Banco Girado} (por ejemplo, Banco Galicia, Santander, Banco Macro, etc.).
    \item Ingresá el \textbf{Número de Cheque} (o el ID de transacción si es un E-Cheq).
    \item En el campo \textbf{Monto}, escribí el valor facial del cheque en pesos.
    \item Indicá la \textbf{Fecha de Emisión} y la \textbf{Fecha de Pago} (el día exacto a partir del cual el cheque se puede depositar o cobrar por ventanilla).
    \item En el campo de archivo adjunto, subí la foto del cheque de papel o el comprobante en PDF del E-Cheq. De esta manera, cualquier compañero de trabajo puede ver el cheque digitalizado en pantalla sin necesidad de ir a buscar el papel físico a la tesorería.
    \item Hacé clic en \textbf{"Guardar Cheque"}. El valor ingresará en estado \texttt{PENDIENTE}.
\end{enumerate}

\textbf{Estados del cheque a lo largo de su vida:}
\begin{itemize}
    \item \textbf{PENDIENTE:} El cheque recién ingresó y está guardado en custodia a la espera de ser cobrado, vendido o debitado.
    \item \textbf{COMPRADO:} Aplica a cheques de terceros que tu empresa le compró a un cliente adelantándole el dinero en mano con descuento de intereses.
    \item \textbf{VENDIDO:} Cheque que tu empresa entregó o descontó ante un banco, inversor o proveedor.
    \item \textbf{ACREDITADO:} Cheque de terceros que ya fue depositado y la plata entró efectivamente en tu cuenta bancaria.
    \item \textbf{DEBITADO:} Cheque propio que se presentó al cobro y el banco ya te restó los fondos de tu cuenta corriente.
\end{itemize}

\textbf{Paso a paso: Cómo comprar o descontar un cheque a un cliente}
Si un cliente o comercio te trae un cheque a cobrar a 30, 60 o 90 días y necesita el dinero en efectivo hoy:
\begin{enumerate}
    \item Ubicá el cheque en la grilla y, en el menú de acciones de la fila, hacé clic en \textbf{"Operar"} y seleccioná \textbf{"Comprar"}.
    \item Seleccioná el cliente o comercio al que le comprás el cheque.
    \item Indicá la \textbf{Fecha de Operación} (la fecha en la que le entregás la plata).
    \item Ingresá la \textbf{Tasa Nominal Anual (TNA) de Descuento} que acordaron cobrarle.
    \item Ingresá el porcentaje de \textbf{Gastos Administrativos} (el sistema te sugiere un 2,80\% habitual, pero podés cambiarlo).
    \item Si el cheque es de otra provincia o el banco demora días hábiles adicionales en acreditarlo, sumá esos días en el campo \textbf{Días de Castigo} (por ejemplo, 1 o 2 días de clearing).
    \item En el recuadro de \textbf{Vista Previa de Liquidación}, vas a ver el cálculo automático al instante:
          \begin{itemize}
              \item Los días corridos totales hasta el vencimiento.
              \item El descuento de intereses financieros.
              \item Los gastos administrativos y el 21\% de IVA correspondiente sobre intereses y gastos.
              \item El \textbf{Importe Neto a Pagar}: la suma exacta de dinero que tenés que transferirle o entregarle en mano al cliente.
          \end{itemize}
    \item Hacé clic en \textbf{"Confirmar Operación"}. El cheque pasa a estado \texttt{COMPRADO} y la ganancia por el descuento queda asentada en el sistema.
\end{enumerate}

\textbf{Paso a paso: Cómo vender o entregar un cheque a un tercero}
Si tenés un cheque en cartera (sea propio o de terceros previamente comprado) y lo querés negociar con un banco o entregarlo en parte de pago a un inversor:
\begin{enumerate}
    \item En la grilla, hacé clic en \textbf{"Operar"} sobre el cheque y elegí \textbf{"Vender"}.
    \item Seleccioná el banco o inversor al que se lo transferís.
    \item Cargá la fecha de entrega y la tasa de descuento acordada.
    \item Al confirmar, el cheque pasará a estado \texttt{VENDIDO} y el sistema actualizará quién es el tenedor actual del documento.
\end{enumerate}

\textbf{Paso a paso: Cómo conciliar cheques con tus cuentas bancarias}
Cuando los cheques se depositan o se debitan en la realidad bancaria, el sistema te permite unirlos con el extracto del banco en dos clics:
\begin{itemize}
    \item \textbf{Para cheques de terceros depositados (Acción Acreditar):} Hacé clic en el botón \textbf{"Acreditar"} en la fila del cheque. Se abrirá una ventana con los movimientos recientes de tu extracto bancario. Elegí el ingreso de dinero correspondiente y presioná \textbf{"Vincular"}. El cheque pasa a estar \texttt{ACREDITADO} y muestra la insignia verde \texttt{[Vinculado]}.
    \item \textbf{Para cheques propios que se pagaron (Acción Debitar):} Hacé clic en el botón \textbf{"Debitar"} en la fila del cheque propio. Buscá el débito en el extracto de tu cuenta corriente y presiona \textbf{"Vincular"}. El cheque pasará a estar \texttt{DEBITADO}.
    \item \textbf{¿Te equivocaste de movimiento?}: Hacé clic sobre la insignia \texttt{[Vinculado]} y el sistema te permite desvincular el cheque de inmediato, devolviéndolo a su estado anterior en cartera sin alterar los registros históricos.
\end{itemize}

\subsection{Inversores y Deuda Privada}

El módulo de \textbf{Inversores y Deuda Privada} está pensado para aquellas empresas que se financian recibiendo fondos de personas o empresas privadas para prestar en la calle. 

Desde acá podés llevar el control milimétrico de cada inversor: quién aportó el dinero, en qué condiciones (plazo y tasa de interés acordada), cuánta renta ganó hasta el día de hoy, y cómo registrar pagos de intereses, rescates o renovaciones.

El módulo se compone de cuatro partes muy sencillas:
\begin{enumerate}
    \item \textbf{Padrón de Inversores:} El registro de las personas o empresas que aportan capital, con sus datos de contacto y cuentas bancarias.
    \item \textbf{Cuentas Comitentes:} Las cuentas donde se radican las inversiones. Pueden pertenecer a una sola persona o a varios cotitulares.
    \item \textbf{Series de Deuda:} Los tramos o pagarés emitidos (por ejemplo, una colocación a 180 días con una tasa fija anual).
    \item \textbf{Movimientos:} El libro diario donde queda anotado cada ingreso de dinero (suscripción), retiro (rescate), pago de intereses o renovación.
\end{enumerate}

\subsubsection{Inversores}
Es el directorio de personas físicas o jurídicas que confían sus fondos a la empresa.

\textbf{Paso a paso: Cómo dar de alta a un nuevo inversor}
\begin{enumerate}
    \item Ingresá a \textbf{Inversores} > \textbf{Padrón de Inversores} y hacé clic en \textbf{"+ Nuevo Inversor"}.
    \item Ingresá su \textbf{CUIT o CUIL} (11 dígitos sin guiones) y su \textbf{Razón Social} o Nombre completo.
    \item El sistema realiza una consulta automática con ARCA (ex AFIP) para verificar que el CUIT exista y que el nombre ingresado coincida con el padrón impositivo oficial.
    \item Completá su \textbf{Domicilio Legal}, su \textbf{Correo Electrónico} (donde recibirá las liquidaciones de intereses) y su \textbf{Teléfono}.
    \item Cargá sus \textbf{Datos Bancarios}: CBU de 22 dígitos, Nombre del Banco y Número de cuenta donde se le transferirán los intereses o devoluciones de capital.
    \item Hacé clic en \textbf{"Guardar Inversor"}.
\end{enumerate}

\subsubsection{Cuentas Comitentes}
Cada inversor necesita tener al menos una cuenta comitente para poder colocar su dinero. 

\textbf{Régimen de Titularidad (¡Muy importante!):}
Cuando una cuenta tiene más de un titular (por ejemplo, dos socios o familiares), podés configurarla bajo dos modalidades bien diferenciadas:
\begin{itemize}
    \item \textbf{A Orden Indistinta (\texttt{conjunta = false}):} Significa que el dinero, los rescates o los intereses pueden ser retirados o cobrados por \textbf{cualquiera de los titulares de forma individual e independiente}, sin necesidad de que los demás titulares intervengan o firmen.
    \item \textbf{A Orden Conjunta (\texttt{conjunta = true}):} Es una cuenta mancomunada. Para disponer o retirar fondos, la operación debe realizarse necesariamente \textbf{a nombre y con la conformidad de todos los titulares juntos}.
\end{itemize}

\textbf{Paso a paso: Cómo abrir una Cuenta Comitente}
\begin{enumerate}
    \item Andá a \textbf{Inversores} > \textbf{Cuentas Comitentes} y hacé clic en \textbf{"+ Nueva Cuenta"}.
    \item Asignale un nombre o identificador (por ejemplo, \texttt{CTA-101}).
    \item Elegí la modalidad de titularidad: \textbf{Indistinta} o \textbf{Conjunta}.
    \item Agregá a los inversores que formarán parte de la cuenta y definí su orden (Titular 1, Titular 2, etc.).
    \item Hacé clic en \textbf{"Guardar Cuenta"}.
    \item En la grilla podés hacer clic en el botón de la lupa o ficha para abrir el \textbf{Estado de Cuenta}, donde vas a ver en cualquier momento el saldo vivo total, los intereses ganados y el historial completo de movimientos.
\end{enumerate}

\subsubsection{Series de Deuda}
Una \textbf{Serie de Deuda} representa el compromiso financiero asumido con los inversores (similar a un pagaré o plazo fijo privado).

\textbf{Paso a paso: Cómo crear una nueva Serie de Deuda}
\begin{enumerate}
    \item Andá a \textbf{Series de Deuda} y hacé clic en \textbf{"+ Nueva Serie"}.
    \item Ingresá un nombre identificatorio (por ejemplo, \texttt{SERIE XXV - 180 DÍAS}).
    \item Indicá la \textbf{Fecha de Suscripción} (cuándo comienza a correr la inversión).
    \item Indicá el \textbf{Plazo Contractual} expresado en días corridos (por ejemplo, 90, 180 o 365 días).
    \item Cargá la \textbf{Tasa Nominal Anual (TNA \%)} que se le pagará al inversor por su dinero.
    \item El sistema calculará automáticamente la \textbf{Fecha de Vencimiento} exacta sumando los días pactados.
    \item Si interviene algún comisionista o socio que acercó al inversor, podés activar la opción de comisiones e indicar el porcentaje correspondiente.
    \item Hacé clic en \textbf{"Guardar Serie"}.
\end{enumerate}

\textbf{¿Cómo calcula el sistema los intereses? (Devengamiento diario sin calculadora)}
En la pantalla de Series de Deuda contás con un selector de \textbf{Fecha de Corte}:
\begin{itemize}
    \item Cada día que pasa, el sistema calcula de forma automática cuántos pesos de interés ganó el inversor hasta esa fecha exacta.
    \item Te muestra el \textbf{Capital Vigente}, los \textbf{Intereses Devengados} hasta hoy, los \textbf{Intereses a Devengar} (lo que falta ganar hasta que venza el plazo) y el total que la empresa deberá desembolsar al vencimiento.
\end{itemize}

\textbf{Paso a paso: Cómo renovar una inversión cuando vence}
Cuando una serie de deuda llega a su fecha de vencimiento:
\begin{enumerate}
    \item En la grilla de Series, hacé clic en el botón \textbf{"Renovar"}.
    \item Elegí la nueva serie receptora donde el inversor colocará sus fondos para el siguiente período.
    \item El asistente te da dos opciones muy claras para los intereses ganados:
          \begin{itemize}
              \item \textbf{Retirar intereses en mano:} La empresa le transfiere o paga en efectivo las ganancias al inversor, y renueva únicamente el capital original aportado.
              \item \textbf{Capitalizar intereses (Reinversión total):} Las ganancias se suman al capital, de modo que la nueva inversión arranca con un monto mayor y genera más intereses hacia adelante.
          \end{itemize}
    \item Hacé clic en \textbf{"Confirmar Renovación"}. El sistema cancela la serie anterior y da de alta la nueva serie con todos sus saldos al día de forma automática.
\end{enumerate}

\subsubsection{Movimientos de Deuda}
Es el libro mayor donde queda registrada cada transacción de dinero con los inversores:
\begin{itemize}
    \item \textbf{Suscripción (+):} Cuando el inversor pone plata en una serie. Aumenta su saldo.
    \item \textbf{Rescate (-):} Cuando el inversor retira total o parcialmente su capital.
    \item \textbf{Retiro de Intereses (-):} Cuando la empresa le paga las ganancias devengadas.
    \item \textbf{Renovaciones (+ / -):} El traspaso automático de fondos de una serie vencida a una nueva.
\end{itemize}

\textbf{Paso a paso: Cómo ingresar una Suscripción de fondos}
\begin{enumerate}
    \item En \textbf{Movimientos}, hacé clic en \textbf{"+ Nuevo Movimiento"}.
    \item Elegí el tipo de movimiento \textbf{Suscripción (+)}.
    \item Seleccioná la Cuenta Comitente del inversor y la Serie de Deuda correspondiente.
    \item Ingresá el monto de dinero depositado y la fecha efectiva de acreditación bancaria.
    \item Hacé clic en \textbf{"Guardar"}. En la grilla vas a ver reflejado el impacto positivo en el saldo de la cuenta comitente y de la serie.
\end{enumerate}

\subsection{Finanzas}

El módulo de \textbf{Finanzas} funciona como la oficina de tesorería y contabilidad de tu empresa: desde acá controlás la plata que tenés en los bancos, registrás los gastos y facturas de proveedores que hay que pagar, liquidás las comisiones de los socios que te traen clientes y calculás cuánto hay que pagarle a ARCA (ex AFIP) por IVA e Ingresos Brutos (IIBB).

El módulo reúne cinco herramientas muy prácticas:
\begin{enumerate}
    \item \textbf{Liquidación de Comisiones:} Para calcular y pagar en segundos lo que le corresponde a los socios comerciales por los préstamos otorgados, las cuotas cobradas o los inversores captados.
    \item \textbf{Bancos y Cuentas:} Para consultar en una sola pantalla los saldos de tus cuentas corrientes, saber cuánta plata tenés rindiendo en Fondos Comunes de Inversión (FCI) o Plazos Fijos, y subir los extractos del banco para conciliar movimientos sin cargarlos a mano.
    \item \textbf{Comprobantes (Cuentas a Pagar):} Para llevar el control de todas las facturas de proveedores (alquiler, luz, papelería, servicios), armar planes de pago en cuotas y registrar las cancelaciones.
    \item \textbf{Posición de IVA:} Para contrastar el débito fiscal de los intereses cobrados en tus cuotas contra el crédito fiscal de tus compras y retenciones bancarias, sabiendo con exactitud cuánto pagar a fin de mes.
    \item \textbf{Posición de IIBB:} Para liquidar mensualmente el Impuesto sobre los Ingresos Brutos provincia por provincia bajo el régimen de Convenio Multilateral, según el domicilio de tus clientes.
\end{enumerate}

\subsubsection{Liquidación de Comisiones}
En este panel calculás de forma automática las retribuciones económicas que le corresponden a tus socios comerciales (los originadores que captan solicitudes y los intermediarios de la red comercial).

\textbf{Filtro mensual de liquidación:}
En la parte superior elegís el año y mes que querés liquidar (por ejemplo, \texttt{2026-10}). El sistema analiza todas las operaciones de ese mes y las agrupa en tres solapas muy claras:
\begin{itemize}
    \item \textbf{Comisiones de Colocación:} Muestra los honorarios generados por el otorgamiento de nuevos créditos durante el mes. Vas a ver el capital total prestado, el porcentaje pactado y cuánto le corresponde a cada parte (\textit{Comisión Originador} y \textit{Comisión Intermediario}).
    \item \textbf{Comisiones de Cobranza:} Calcula las comisiones sobre las cuotas que tus clientes pagaron efectivamente en el mes, evitando pagar comisiones por adelantado sobre dinero aún no cobrado.
    \item \textbf{Comisiones de Inversores:} Para las empresas que trabajan con deuda privada, liquida los honorarios pactados con aquellos intermediarios que acercaron fondos de inversores a las Series de Deuda.
\end{itemize}

\textbf{Paso a paso: Cómo liquidar las comisiones del mes y exportarlas a Excel}
\begin{enumerate}
    \item Ingresá a \textbf{Finanzas} > \textbf{Liquidación de Comisiones}.
    \item En el selector superior, elegí el \textbf{Mes de Liquidación} que deseás procesar.
    \item Recorré las solapas (\textbf{Colocación}, \textbf{Cobranza} e \textbf{Inversores}) para verificar las operaciones computadas y los porcentajes aplicados a cada socio.
    \item Si tu empresa le realizó adelantos de dinero al socio durante el mes, el sistema deduce automáticamente esos importes del saldo a cobrar.
    \item En la barra inferior vas a ver el resumen consolidado con el total a pagar. Hacé clic en el botón verde \textbf{"Exportar a Excel"} para descargar la planilla detallada, lista para respaldar la orden de pago bancaria o solicitarle la factura al socio.
\end{enumerate}

\subsubsection{Bancos y Cuentas}
Esta pantalla es el centro de control de tesorería: te permite supervisar la liquidez de tu empresa y mantener tus cuentas bancarias conciliadas al día.

\textbf{Panel de Control de Fondos (Dashboard de Tesorería):}
Te muestra en tiempo real y para cualquier rango de fechas:
\begin{itemize}
    \item \textbf{Saldo Disponible:} La plata líquida que tenés en tus cuentas bancarias para operar en el día a día.
    \item \textbf{Saldo en FCI:} Dinero colocado transitoriamente en Fondos Comunes de Inversión de rescate inmediato (Money Market) para que no pierda valor contra la inflación.
    \item \textbf{Saldo en Plazo Fijo:} Colocaciones bancarias a término devengando intereses.
    \item \textbf{Ingresos y Egresos del período:} El total de dinero que entró (por cobranzas, rescates o transferencias) y todo lo que salió (por préstamos otorgados, gastos o pago a proveedores).
    \item \textbf{Flujo Neto:} La diferencia final ($\text{Ingresos} - \text{Egresos}$), que te dice si el período cerró en superávit o déficit de fondos.
\end{itemize}

\textbf{Subida de extractos bancarios y conciliación automática:}
Para que no tengas que escribir a mano cada movimiento del extracto bancario, el sistema puede leer directamente los resúmenes descargados de los principales bancos (Banco BICA, Banco Santander, Banco de La Pampa, Banco Galicia, Banco Macro, etc.). El sistema analiza cada renglón y reconoce automáticamente conceptos habituales (como transferencias, débitos impositivos, comisiones de mantenimiento o pagos de VEP).

\textbf{Paso a paso: Cómo cargar y conciliar un extracto bancario}
\begin{enumerate}
    \item Ingresá a \textbf{Finanzas} > \textbf{Bancos y Cuentas} y seleccioná la cuenta bancaria sobre la que querés trabajar.
    \item Hacé clic en el botón \textbf{"Subir Extracto"}.
    \item Seleccioná el archivo de extracto descargado de tu home banking (en formato texto o Excel).
    \item El sistema procesará el archivo y te mostrará la grilla con todos los renglones importados.
    \item Los movimientos reconocidos quedarán clasificados de inmediato. Si tenés varios movimientos similares sin catalogar (por ejemplo, comisiones bancarias), podés tildarlos juntos y hacer clic en \textbf{"Asignar Concepto en Bloque"} para clasificarlos en un solo paso.
    \item Si el extracto incluye el cobro o débito de un cheque en cartera, hacé clic en el botón de vinculación para unirlo directamente con el cheque correspondiente.
\end{enumerate}

\subsubsection{Comprobantes (Cuentas a Pagar a Proveedores)}
En este módulo administrás las compras y gastos de la empresa (alquileres, luz, papelería, servicios contables, software, etc.), asegurando un control ordenado de tus deudas comerciales y los vencimientos de pago.

\textbf{Padrón de Proveedores:}
Registrás los datos comerciales de tus acreedores: CUIT, Razón Social, Condición frente al IVA, domicilio, canales de contacto y el CBU bancario donde se les transfiere.

\textbf{Carga y control de comprobantes:}
\begin{itemize}
    \item \textbf{Tipos de comprobantes admitidos:} Facturas A, B, C, M, Notas de Débito, Notas de Crédito y Recibos.
    \item \textbf{Desglose impositivo completo:} Al cargar una factura, ingresás el importe neto gravado, la alícuota de IVA correspondiente (21\%, 10,5\% o 27\%), las percepciones provinciales o nacionales y los conceptos no gravados. Esto permite que el sistema alimente automáticamente tu liquidación impositiva mensual sin que tengas que volver a cargar los datos.
    \item \textbf{Estados del comprobante:}
          \begin{itemize}
              \item \textit{Pendiente:} Factura cargada que aún no fue pagada.
              \item \textit{Pagado:} Comprobante cancelado en su totalidad.
              \item \textit{Vencido:} Comprobante impago cuya fecha límite de pago ya expiró.
          \end{itemize}
\end{itemize}

\textbf{Paso a paso: Cómo cargar una factura de proveedor y registrar su pago}
\begin{enumerate}
    \item En el menú lateral, andá a \textbf{Finanzas} > \textbf{Comprobantes} y hacé clic en \textbf{"+ Nuevo Comprobante"}.
    \item Seleccioná el \textbf{Proveedor} emisor (si no existe, podés crearlo en el momento).
    \item Elegí el \textbf{Tipo de Comprobante} (por ejemplo, Factura A), indicá el punto de venta y el número de comprobante.
    \item Ingresá la \textbf{Fecha de Emisión} y la \textbf{Fecha de Vencimiento} para el pago.
    \item Cargá el importe neto gravado y seleccioná la alícuota de IVA. Si la factura incluye percepciones de Ingresos Brutos o IVA, sumalas en sus respectivos casilleros. El sistema calculará el total general automáticamente.
    \item Hacé clic en \textbf{"Guardar Comprobante"}. La factura quedará asentada en estado \textit{Pendiente}.
    \item Cuando realices el pago (por transferencia bancaria o entregando un cheque de cartera), hacé clic en \textbf{"Pagar / Cancelar"} en la fila del comprobante, seleccioná el medio de pago utilizado y confirmá. La factura pasará a estado \textit{Pagado}.
\end{enumerate}

\subsubsection{Posición de IVA}
Esta herramienta le ahorra horas de trabajo a tu área contable, calculando de manera automática la declaración jurada mensual del Impuesto al Valor Agregado:

\begin{itemize}
    \item \textbf{Débito Fiscal (IVA Ventas):} Es el IVA acumulado por los intereses cobrados y facturados a los clientes en sus cuotas durante el mes.
    \item \textbf{Crédito Fiscal (IVA Compras):} Es el IVA computable de todas las facturas de proveedores que cargaste en el módulo de Comprobantes durante el mismo período.
    \item \textbf{Retenciones Bancarias Sufridas:} El dinero que los bancos te retuvieron en concepto de IVA en tus cuentas durante el mes.
    \item \textbf{Percepciones de Compras:} Percepciones de IVA liquidadas por proveedores en sus facturas de compras.
    \item \textbf{Saldo a Favor del Período Anterior:} Crédito fiscal acumulado que te quedó a favor del mes pasado.
    \item \textbf{Pagos a Cuenta (VEPs):} Anticipos o volantes de pago ya abonados ante la agencia tributaria.
\end{itemize}

\textbf{Cálculo de la Posición Neta:}
El sistema aplica la fórmula oficial de liquidación:
\begin{equation}
    \begin{aligned}
        \text{Posición Neta IVA} ={} & \text{IVA Ventas} - \text{IVA Compras} \\
        & - \text{Retenciones} - \text{Percepciones} + \text{Saldo Anterior} - \text{Pagos VEP}
    \end{aligned}
\end{equation}
Si el resultado es positivo, representa el importe que la empresa debe pagarle a ARCA/AFIP; si es negativo, constituye un saldo técnico a favor que se traslada automáticamente al mes siguiente.

\textbf{Paso a paso: Cómo revisar y guardar la Posición mensual de IVA}
\begin{enumerate}
    \item Ingresá a \textbf{Finanzas} > \textbf{Posición de IVA}.
    \item Seleccioná el \textbf{Año y Mes} que deseás liquidar (por ejemplo, \texttt{2026-10}).
    \item Verificá las cifras automáticas que el sistema calculó a partir de tus ventas y compras del mes.
    \item Si tu contador te indica algún ajuste especial o querés ingresar un pago a cuenta, cargalo en el campo correspondiente con su justificación.
    \item Hacé clic en \textbf{"Guardar Declaración"}. La liquidación quedará archivada en el historial inmutable de la empresa y podés exportar el reporte a Excel para respaldar la presentación impositiva oficial.
\end{enumerate}

\subsubsection{Posición de IIBB (Ingresos Brutos / Convenio Multilateral)}
Este módulo calcula mensualmente el Impuesto sobre los Ingresos Brutos distribuyendo la base imponible según la provincia donde vive cada cliente, tal como lo exige el régimen de Convenio Multilateral.

\textbf{Cálculo territorial automático:}
\begin{itemize}
    \item \textbf{Base Gravada por Provincia:} El sistema analiza cada cuota cobrada en el mes y, según el domicilio legal del titular del crédito, imputa los intereses financieros netos a la provincia que corresponde.
    \item \textbf{Alícuotas Provinciales Configurables:} Te permite establecer y actualizar el porcentaje impositivo vigente para cada jurisdicción (por ejemplo, 5\%, 7\%, 8\%).
    \item \textbf{Impuesto Determinado por Jurisdicción:}
          \begin{equation}
              \text{Impuesto Provincial}_{j} = \text{Base Gravada}_{j} \times \left(\frac{\text{Al\'{\i}cuota}_{j}}{100}\right)
          \end{equation}
          donde $j$ representa a cada una de las provincias donde tenés operaciones vigentes.
\end{itemize}

\textbf{Determinación del Saldo a Pagar:}
El sistema consolida el impuesto total de todas las provincias y descuenta las retenciones bancarias provinciales, percepciones sufridas y pagos a cuenta:
\begin{equation}
    \begin{aligned}
        \text{Saldo IIBB} ={} & \sum_{j} \text{Impuesto Provincial}_{j} - \text{Retenciones} \\
        & - \text{Percepciones} + \text{Saldo Anterior} - \text{Pagos VEP}
    \end{aligned}
\end{equation}

\textbf{Paso a paso: Cómo liquidar y presentar la Posición de IIBB}
\begin{enumerate}
    \item Ingresá a \textbf{Finanzas} > \textbf{Posición de IIBB}.
    \item Elegí el período mensual a liquidar.
    \item Vas a ver en pantalla la tabla con todas las provincias alcanzadas: el monto de base gravada, la alícuota aplicable y el impuesto calculado para cada una de ellas.
    \item Revisá el total de retenciones y percepciones provinciales acumuladas.
    \item Cambiá el estado de la liquidación a \textbf{"Presentada"} y, una vez que abones el volante electrónico de pago, marcala como \textbf{"Pagada"} registrando la fecha de pago.
    \item Descargá la planilla resumen en Excel para tu archivo y para el profesional contable.
\end{enumerate}

\subsection{Configuración}

El módulo de \textbf{Configuración} es el centro de control de Credit Manager. Su principal ventaja es que le otorga a tu empresa \textbf{autonomía total}: podés modificar tasas de interés, crear nuevos esquemas de comisiones, ajustar los límites de riesgo crediticio, dar de alta usuarios y personalizar los modelos de contratos en Word sin depender de programadores ni requerir cambios técnicos en el sistema.

El módulo se compone de cuatro áreas:
\begin{enumerate}
    \item \textbf{Tablas Auxiliares:} El catálogo maestro con las 18 tablas paramétricas del sistema (provincias, socios, tasas, empleadores, bancos, etc.).
    \item \textbf{Acciones del Sistema:} Herramientas de mantenimiento general para sincronizar estados de cobro, activar o apagar módulos y personalizar el logo corporativo.
    \item \textbf{Usuarios:} Administración de operadores, asignación de permisos según el puesto de trabajo y restablecimiento de contraseñas.
    \item \textbf{Papelería:} Carga y personalización de plantillas en Word (\texttt{.docx}) con reemplazo automático de variables para generar contratos y pagarés en segundos.
\end{enumerate}

\subsubsection{Tablas Auxiliares}
En esta sección consultás y modificás las 18 tablas que rigen el funcionamiento del negocio. En cada tabla podés filtrar por cualquier columna, ordenar los registros de menor a mayor, exportar la información a Microsoft Excel y crear o editar renglones con un simple formulario modal.

Las 18 tablas se dividen en seis áreas de trabajo:

\begin{enumerate}
    \item \textbf{Geografía y Normativa:}
          \begin{itemize}
              \item \textbf{Nacionalidades (\texttt{nacionalidades}):} Listado de países. Incluye una marca de alerta para países de alto riesgo según el GAFI (Grupo de Acción Financiera Internacional), advirtiendo si corresponde aplicar una debida diligencia reforzada en Prevención de Lavado (PLA/FT).
              \item \textbf{Provincias (\texttt{provincias}):} Catálogo de provincias argentinas, fundamental para calcular las tasas del Impuesto de Sellos y la distribución de Ingresos Brutos (Convenio Multilateral).
          \end{itemize}

    \item \textbf{Socios Comerciales y Red de Ventas:}
          \begin{itemize}\raggedright
              \item \textbf{Socios Comerciales (\texttt{socios}):} Empresas y mutuales que actúan como originadores de créditos o intermediarios. Acá se configuran sus datos de contacto, CUIT, cuenta bancaria para transferirles comisiones, día de corte mensual y los \textit{Anticipos Vigentes} (adelantos de dinero realizados para compensar en liquidaciones futuras).
              \item \textbf{Empleadores (\texttt{empleadores}):} Organismos pagadores y empresas sobre cuyos recibos de sueldo se aplican los descuentos de cuotas. La casilla \textit{Es Pasivo} diferencia entre cajas de previsión social (jubilados y pensionados) y personal en actividad.
              \item \textbf{Comercializadores (\texttt{comercializadores}):} Vendedores y asesores comerciales externos que acercan solicitudes de crédito.
              \item \textbf{Relaciones Mapeadas (\texttt{relaciones}):} Diccionario de equivalencias que permite vincular los códigos de clientes o créditos de un socio externo con los identificadores internos de Credit Manager durante la importación masiva de carteras.
          \end{itemize}

    \item \textbf{Operadores de Cheques:}
          \begin{itemize}
              \item \textbf{Operadores (\texttt{operadores}):} Directorio de personas y entidades con las que se negocian cheques, con su calificación crediticia (\textit{Excelente}, \textit{Bueno}, \textit{Regular}, etc.) para controlar límites y aforos de descuento en tesorería.
          \end{itemize}

    \item \textbf{Tasas y Comisiones:}
          \begin{itemize}\raggedright
              \item \textbf{Tasas y Comisiones (\texttt{tasas\_y\_comisiones}):} Las condiciones financieras aplicables a los créditos: plazos en cuotas, Tasa Nominal Anual con IVA (\texttt{tna\_c\_iva}), comisiones para originadores e intermediarios, gastos administrativos de otorgamiento, porcentaje de sellado provincial y el factor de recargo punitorio por mora.
              \item \textbf{Comisiones de Deuda (\texttt{comisiones\_deuda}):} Porcentajes de comisión acordados para socios que acercan inversores a las Series de Deuda.
          \end{itemize}

    \item \textbf{Finanzas y Tesorería:}
          \begin{itemize}
              \item \textbf{Bancos (\texttt{bancos}):} Entidades bancarias del sistema financiero con el formato de extracto compatible configurado para conciliar automáticamente.
              \item \textbf{Cuentas Bancarias (\texttt{cuentas}):} Cuentas bancarias propias con su CBU, Alias, número de cuenta y moneda.
              \item \textbf{Conceptos (\texttt{conceptos}):} Motivos de ingresos y egresos de dinero (por ejemplo, cobranzas, sueldos, alquiler, débitos bancarios).
              \item \textbf{Clasificaciones (\texttt{clasificaciones}):} Agrupaciones presupuestarias para los informes gerenciales de flujo de fondos.
          \end{itemize}

    \item \textbf{Riesgo Crediticio y Capacidad de Pago (PLA/FT):}
          \begin{itemize}\raggedright
              \item \textbf{Factores de Riesgo (\texttt{factores\_riesgo}):} Criterios evaluados en la matriz de scoring crediticio, con sus pesos y ponderaciones.
              \item \textbf{Multiplicadores de Riesgo (\texttt{multiplicadores\_riesgo}):} Coeficientes de ajuste según las características del tomador.
              \item \textbf{Perfiles Transaccionales (\texttt{reglas\_perfiles\_transaccionales}):} Reglas para calcular con total precisión la capacidad de pago del solicitante a partir de su recibo de haberes: define qué conceptos suman (básico, adicionales, horas extras, asignaciones) y qué descuentos restan, determinando el \textbf{Sueldo Computable} y la \textbf{Cuota Afectable Máxima}:
                    \begin{equation}
                        \text{Sueldo Computable} = \text{Básico} + \sum_{k \in \text{Habilitados}} \text{Concepto}_{k} - \sum_{m \in \text{Deducciones}} \text{Descuento}_{m}
                    \end{equation}
                    \begin{equation}
                        \text{Cuota Afectable Máxima} = \text{Sueldo Computable} \times \text{Cupo (\%)}
                    \end{equation}
              \item \textbf{Punitorios (\texttt{penalties}):} Escalas de recargo moratorio según los días de atraso.
              \item \textbf{Políticas Crediticias (\texttt{politicas\_crediticias}):} Requisitos de admisión por socio originador: edades mínimas y máximas de otorgamiento (para hombres, mujeres y jubilados) y antigüedad laboral mínima requerida.
          \end{itemize}
\end{enumerate}

\textbf{Paso a paso: Cómo consultar o editar parámetros en las Tablas Auxiliares}
\begin{enumerate}
    \item En el menú lateral, ingresá a \textbf{Configuración} > \textbf{Tablas Auxiliares}.
    \item En el selector desplegable superior, elegí la tabla que necesitás consultar (por ejemplo, \textit{Tasas y Comisiones} o \textit{Empleadores}).
    \item Para modificar un registro existente, hacé clic en el botón de edición (ícono de lápiz) en la fila correspondiente. Se abrirá una ventana donde podés actualizar valores (como la tasa TNA o un teléfono de contacto).
    \item Para dar de alta un nuevo registro, hacé clic en el botón verde superior \textbf{"+ Nuevo Registro"}, completá los campos del formulario y presioná \textbf{"Guardar"}.
    \item El cambio surtirá efecto inmediato en todo el sistema sin necesidad de reiniciar la aplicación.
\end{enumerate}

\subsubsection{Acciones del Sistema}
Este panel agrupa tareas de mantenimiento general, personalización visual y supervisión de seguridad:

\begin{itemize}
    \item \textbf{Sincronización Global de Estados:}
          Es el botón de puesta al día general. Al ejecutarlo, el sistema revisa en pocos segundos el universo completo de créditos, cuotas y clientes, comparando la fecha actual contra los vencimientos y pagos asentados. Actualiza en tiempo real los estados de mora (\textit{Al Día}, \textit{Mora Leve}, \textit{Mora Media}, \textit{Mora Grave}, \textit{Incobrable} o \textit{Cancelado}), asegurando que las cifras de los reportes y del Dashboard sean exactas.

    \item \textbf{Activar o Desactivar Módulos (\textit{Feature Toggling}):}
          Panel exclusivo para el rol de \textbf{Administrador}. Si tu empresa actualmente no opera con cheques o inversores privados, podés apagar esos módulos con un simple interruptor para que no aparezcan en el menú lateral. Cuando decidas utilizarlos, los activás con un clic y quedarán disponibles al instante para todos los usuarios.

    \item \textbf{Identidad Corporativa (Logo y Favicon):}
          Permite subir la imagen del logo de tu empresa (para que aparezca en el encabezado de la plataforma y en los contratos o recibos emitidos) y el ícono de la pestaña del navegador (\textit{favicon}).

    \item \textbf{Registro de Auditoría (Bitácora de Seguridad):}
          El sistema registra de forma automática e inalterable cada acción importante que se realiza: qué usuario la ejecutó, a qué hora exacta, desde qué computadora y qué datos modificó. Esto garantiza total transparencia operativa y brinda respaldo ante cualquier consulta interna.
\end{itemize}

\textbf{Paso a paso: Cómo ejecutar la Sincronización Global de Estados}
\begin{enumerate}
    \item Ingresá a \textbf{Configuración} > \textbf{Acciones del Sistema}.
    \item En la tarjeta de \textbf{Sincronización de Estados}, hacé clic en el botón azul \textbf{"Sincronizar Estados Ahora"}.
    \item El sistema procesará todas las cuentas y mostrará un mensaje de confirmación indicando la cantidad de créditos y cuotas verificadas.
\end{enumerate}

\subsubsection{Usuarios y Permisos}
Para organizar el trabajo en equipo y proteger la confidencialidad de la información, el sistema asigna perfiles de acceso (roles) adaptados a cada tarea en la empresa:

\begin{enumerate}
    \item \textbf{Administrador:} Acceso completo a todas las pantallas, configuración global de parámetros, gestión de usuarios y auditoría.
    \item \textbf{Gerente:} Supervisión ejecutiva. Consulta paneles consolidados, reportes patrimoniales y autoriza excepciones en políticas crediticias.
    \item \textbf{Oficial de Crédito:} Para el equipo comercial: alta y consulta de clientes, simulación de cuotas, evaluación de perfiles transaccionales y otorgamiento de préstamos.
    \item \textbf{Operador de Cobranzas:} Para el equipo de recupero: registro de cobros individuales, procesamiento de planillas masivas de Excel y seguimiento de atrasos.
    \item \textbf{Operador de Inversiones:} Para el área de fondeo: administración de cuentas comitentes, series de deuda, depósitos y renovaciones.
    \item \textbf{Responsable de Finanzas:} Para la tesorería: conciliación de extractos bancarios, carga de facturas de proveedores y liquidación de comisiones, IVA e IIBB.
    \item \textbf{Auditor / Solo Lectura:} Para revisores o contadores externos: consulta de información y descarga de reportes sin posibilidad de modificar ni borrar datos.
\end{enumerate}

\textbf{Paso a paso: Cómo dar de alta un usuario o cambiar una contraseña}
\begin{enumerate}
    \item Andá a \textbf{Configuración} > \textbf{Usuarios}.
    \item Para crear un operador nuevo, hacé clic en el botón verde \textbf{"+ Nuevo Usuario"}. Escribí su nombre completo, su correo electrónico (con el que iniciará sesión) y asignale el rol correspondiente.
    \item Para cambiar la clave de un usuario que la olvidó, hacé clic en el botón de la llavecita (\textit{Cambiar Contraseña}) en su fila, ingresá la nueva clave y confirmá.
    \item Si un colaborador deja de trabajar en la empresa, simplemente editá su usuario y desmarcá la casilla \textbf{Activo}. Su acceso quedará bloqueado de inmediato, pero todas sus operaciones históricas quedarán debidamente preservadas en la auditoría.
\end{enumerate}

\subsubsection{Papelería (Plantillas Word y Emisión de Documentos)}
Este módulo automatiza por completo la confección de contratos, pagarés y solicitudes legales, eliminando la tarea lenta y propensa a errores de rellenar documentos a mano en Word.

\textbf{¿Cómo funciona la sustitución de variables?}
Tu empresa sube sus propios modelos de contratos en formato Word (\texttt{.docx}). En las partes del texto donde van los datos del cliente, del préstamo o de la empresa, se colocan etiquetas entre llaves dobles, por ejemplo: \texttt{\{\{ cliente.nombre \}\}} o \texttt{\{\{ credito.monto\_otorgado\_letras \}\}}.
Al momento de otorgar un crédito o vender una cartera, el sistema toma la plantilla y reemplaza al instante cada etiqueta por el dato real, formateando los importes en números y letras automáticamente.

\textbf{Categorías de documentación:}
\begin{itemize}
    \item \textbf{Créditos:} Solicitud de Préstamo, Pagaré a la Vista con Cláusula sin Protesto, Autorización de Débito Bancario o Descuento de Haberes, Declaración Jurada de Origen y Licitud de Fondos (PLA/FT), Condición de Persona Expuesta Políticamente (PEP), Constancia de Consulta en RePET y Seguro de Vida sobre Saldo Deudor.
    \item \textbf{Ventas de Cartera:} Contrato Marco de Cesión de Créditos, Anexos de Cartera (con o sin recurso) y Notificaciones de cesión a los deudores.
\end{itemize}

\textbf{Plantillas generales o por socio comercial:}
Podés configurar documentos de uso general compartidos por toda la organización o definir plantillas exclusivas para un socio comercial determinado (respetando su membrete, domicilios especiales y cláusulas contractuales propias).

\textbf{Variables disponibles para utilizar en tus plantillas Word:}
\begin{itemize}\raggedright
    \item \textbf{Datos del Cliente:} \texttt{cliente.nombre}, \texttt{cliente.apellido}, \texttt{cliente.dni}, \texttt{cliente.cuil}, \texttt{cliente.domicilio}, \texttt{cliente.calle}, \texttt{cliente.calle\_nro}, \texttt{cliente.piso}, \texttt{cliente.depto}, \texttt{cliente.cp}, \texttt{cliente.cbu}, \texttt{cliente.cuenta\_bancaria}, \texttt{cliente.banco}, \texttt{cliente.localidad}, \texttt{cliente.provincia}, \texttt{cliente.nacionalidad}, \texttt{cliente.estado\_civil}, \texttt{cliente.telefono}, \texttt{cliente.mail}, \texttt{cliente.sexo}, \texttt{cliente.fecha\_nacimiento}, \texttt{cliente.cargo}, \texttt{cliente.pep}, \texttt{cliente.repet}, \texttt{cliente.repet\_id}, \texttt{cliente.repet\_fecha\_consulta}, \texttt{cliente.repet\_estado}.
    \item \textbf{Referidos Personales:} \texttt{referido\_1.nombre}, \texttt{referido\_1.apellido}, \texttt{referido\_1.telefono}, \texttt{referido\_1.email}, \texttt{referido\_2.nombre}, \texttt{referido\_2.apellido}, \texttt{referido\_2.telefono}, \texttt{referido\_2.email}.
    \item \textbf{Datos del Empleador:} \texttt{empleador.razon\_social}, \texttt{empleador.fecha\_ingreso}, \texttt{empleador.ingreso\_mensual}, \texttt{empleador.legajo}, \texttt{empleador.domicilio\_calle}, \texttt{empleador.domicilio\_nro}, \texttt{empleador.domicilio\_piso}, \texttt{empleador.domicilio\_depto}, \texttt{empleador.localidad}, \texttt{empleador.provincia}, \texttt{empleador.codigo\_postal}, \texttt{empleador.telefono}.
    \item \textbf{Condiciones Financieras del Crédito:} \texttt{credito.monto\_otorgado}, \texttt{credito.monto\_otorgado\_letras}, \texttt{credito.monto\_neto}, \texttt{credito.monto\_neto\_letras}, \texttt{credito.gastos\_otorgamiento}, \texttt{credito.gastos\_otorgamiento\_letras}, \texttt{credito.sellado}, \texttt{credito.sellado\_letras}, \texttt{credito.interes}, \texttt{credito.interes\_letras}, \texttt{credito.iva}, \texttt{credito.iva\_letras}, \texttt{credito.plazo}, \texttt{credito.valor\_cuota}, \texttt{credito.monto\_total}, \texttt{credito.monto\_total\_letras}, \texttt{credito.tna\_c\_iva}, \texttt{credito.tna\_s\_iva}, \texttt{credito.tea\_s\_iva}, \texttt{credito.tem\_s\_iva}, \texttt{credito.cft}, \texttt{credito.cuotas}, \texttt{credito.fecha\_emision\_dia}, \texttt{credito.fecha\_emision\_mes\_letras}, \texttt{credito.fecha\_emision\_anio}, \texttt{credito.primer\_vto}, \texttt{credito.tabla\_transferencias}.
    \item \textbf{Datos de la Empresa y Socio:} \texttt{empresa.razon\_social}, \texttt{empresa.cuit}, \texttt{empresa.cbu}, \texttt{empresa.nro\_cta\_bancaria}, \texttt{empresa.domicilio}, \texttt{socio.razon\_social}.
    \item \textbf{Cesión de Cartera:} \texttt{cartera.id}, \texttt{cartera.nombre}, \texttt{cartera.fecha\_compra}, \texttt{cartera.monto\_total}, \texttt{cartera.monto\_total\_letras}, \texttt{cartera.socio\_comercial}, \texttt{cartera.recurso}, \texttt{cartera.fecha\_dia}, \texttt{cartera.fecha\_mes\_letras}, \texttt{cartera.fecha\_anio}.
\end{itemize}

\textbf{Paso a paso: Cómo subir una plantilla de Word y emitir la documentación de un crédito}
\begin{enumerate}
    \item En \textbf{Configuración} > \textbf{Papelería}, hacé clic en \textbf{"+ Nueva Plantilla"}.
    \item Asignale un nombre descriptivo (por ejemplo, \textit{Pagaré con Cláusula sin Protesto}), seleccioná la categoría (\textit{Créditos} o \textit{Venta de Cartera}) y adjuntá el archivo Word (\texttt{.docx}) con las etiquetas preparadas.
    \item Establecé el número de orden en que querés que aparezca en el legajo y hacé clic en \textbf{"Guardar"}.
    \item Cuando otorgues un crédito o consultes la ficha de un préstamo existente, hacé clic en el botón \textbf{"Descargar Papelería"}. Vas a poder descargar cada documento en formato Word individual o compilar en el acto un único archivo \textbf{PDF integrado} con todos los formularios completos y ordenados, listo para imprimir y hacer firmar al cliente.
\end{enumerate}

\section{Soporte y Pruebas}
Para cualquier duda operativa sobre el uso de la plataforma o consultas sobre nuevas configuraciones, ponete en contacto con el equipo de soporte interno o el administrador de sistemas de tu empresa. El sistema incluye además un entorno de simulación para que los supervisores puedan poner a prueba nuevas reglas de negocio o cargas masivas de prueba en un ámbito seguro antes de aplicarlas a la operatoria diaria.

\end{document}
